\documentclass[12pt,a4paper]{article}
\usepackage[utf8]{inputenc}
\usepackage[T1]{fontenc}
\usepackage[spanish]{babel}
\usepackage[margin=2.5cm]{geometry}
\usepackage{graphicx}
\usepackage{booktabs}
\usepackage{csquotes}
\usepackage[hidelinks]{hyperref}

% Activar solo cuando referencias.bib contenga fuentes reales verificadas:
% \usepackage[backend=biber,style=ieee]{biblatex}
% \addbibresource{referencias.bib}

\newcommand{\pendiente}[1]{\par\noindent\textit{Pendiente: #1}\par}

\begin{document}

\begin{titlepage}
\centering
{\Large Universidad: pendiente de completar\par}
\vspace{0.5cm}
Facultad: pendiente de completar\par
Carrera: pendiente de completar\par
Asignatura: Sistemas Distribuidos\par
\vfill
{\LARGE\bfseries TA-IND-01\par}
\vspace{0.5cm}
{\Large Informe Académico Individual\par}
\vspace{0.5cm}
{\Large Análisis crítico de SeaweedFS y comparación con MinIO\par}
\vfill
Estudiante: pendiente de completar\par
Paralelo: pendiente de completar\par
Fecha de entrega: 15 de noviembre de 2026\par
\vspace{0.5cm}
\url{https://github.com/jcedenoi4/TA-IND-01}\par
\end{titlepage}

% Plantilla inicial: no constituye el informe final ni cumple aún la extensión.
\section*{Resumen}
\pendiente{Redactar entre 150 y 200 palabras después del análisis y del experimento; incluir enfoque y conclusiones reales.}

\section{Introducción}
\pendiente{Explicar qué problema resuelve SeaweedFS, quién lo utiliza, a qué escala y por qué se eligió. Respaldar las afirmaciones con fuentes.}

\section{Arquitectura del sistema}
\pendiente{Describir los componentes y sus comunicaciones según la versión estudiada. Incorporar un diagrama original de draw.io o TikZ y explicar la evolución desde el esquema a mano del Anexo A.}

\section{Transparencias de distribución}
\pendiente{Analizar acceso, ubicación, migración, reubicación, replicación, concurrencia y fallos; justificar cuáles ofrece y cuáles limita.}

\section{Teorema CAP y PACELC}
\pendiente{Argumentar el comportamiento durante particiones y el compromiso entre latencia y consistencia sin particiones. Precisar operación, configuración y evidencia; evitar clasificaciones sin respaldo.}

\section{Modelos de comunicación}
\pendiente{Identificar cliente-servidor, peer-to-peer o modelo híbrido, sincronía y asincronía, y uso de sockets, RPC o mensajería según fuentes verificadas.}

\section{Tolerancia a fallos}
\pendiente{Explicar fallos soportados, detección y recuperación. Incorporar un incidente real documentado si existe y se verifica.}

\section{Experimentación propia}
\pendiente{Documentar entorno, versiones, hipótesis, procedimiento, comandos, salidas reales e interpretación. Incluir al menos cuatro capturas propias con datos identificativos y fecha.}

\section{Ejemplos del entorno del estudiante}
\subsection{Primer ejemplo}
\pendiente{Describir un sistema conocido de primera mano, una observación real y su relación con un concepto de la unidad.}
\subsection{Segundo ejemplo}
\pendiente{Describir otro sistema y una observación real; distinguir lo observado de las inferencias sobre su implementación.}

\section{Análisis comparativo: SeaweedFS y MinIO}
\pendiente{Construir una tabla con al menos cinco criterios de la unidad, especificando versiones, configuración y fuentes; discutir los compromisos observados.}

\section{Postura crítica}
\pendiente{Argumentar fortalezas, limitaciones y contextos en los que SeaweedFS no sería la mejor opción, con respaldo documental y experimental.}

\section{Conclusiones}
\pendiente{Formular conclusiones concretas derivadas del análisis y de resultados efectivamente obtenidos.}

\section{Declaración de uso de inteligencia artificial}
Utilicé Codex para la revisión inicial del proyecto y la preparación de la estructura de carpetas, el README, la bitácora, el archivo de exclusiones de Git y la plantilla de LaTeX. En esta etapa no se generaron referencias académicas, evidencias ni resultados experimentales.
\pendiente{Antes de entregar, completar el modelo o versión, los usos posteriores y los procedimientos de verificación efectivamente realizados. La declaración debe reflejar el uso total de IA durante el proyecto.}

% Habilitar con biblatex y las referencias verificadas:
% \printbibliography[title={Referencias}]
\section*{Referencias}
\pendiente{Incorporar mediante biblatex al menos seis fuentes de los últimos cinco años, incluidas tres de revistas indexadas con DOI comprobado y metadatos coincidentes.}

\appendix
\section{Esquema inicial a mano}
\pendiente{Incluir la fotografía original del dibujo realizado por el estudiante, con su nombre y fecha visibles.}

\section{Lectura de las fuentes}
\pendiente{Incluir al menos tres capturas de artículos con anotaciones personales, número de página y correspondencia con citas del informe.}

\end{document}
